\begin{abstract}
Modern Large Language Model (LLM) inference workloads
    are increasingly bottlenecked by GPU memory capacity:
    flagship model sizes grow roughly $410\times$ every two years,
    while GPU memory only doubles in the same period.
To run these models on a single GPU,
    today's serving frameworks resort to bespoke,
    application-level offloading of model weights and KV caches
    to host memory~\cite{flexgen, infinigen, headinfer, vllm}.
Each such framework re-implements its own memory management,
    is brittle to model and workload changes,
    and provides little benefit to the broader ecosystem.

Unified Virtual Memory (UVM) is, in principle,
    a natural systems-level answer to this problem:
    it lets the GPU transparently access CPU and peer-GPU memory,
    enabling oversubscription without any application-level offloading code.
In practice, UVM has been considered too slow to use in production:
    PyTorch's own documentation advises against \texttt{cudaMallocManaged},
    and prior work reports that UVM is over $5\times$ slower than
    application-level offloading on LLM workloads~\cite{infinigen}.
The root cause is two-fold:
    (i)~UVM lacks any knowledge of application-level access patterns,
        so it incurs long page-fault and migration latencies; and
    (ii)~modern AI frameworks such as PyTorch and vLLM
        provide no first-class support for UVM,
        so UVM bypasses their caching and allocation optimizations.

We present {\uvmname}, a framework-aware UVM substrate for PyTorch
    that closes this performance gap and turns UVM
    into a practical tool for LLM inference.
{\uvmname} introduces three pieces:
    a \emph{custom UVM-aware caching allocator}
        that actively returns cached memory to CUDA when the GPU is oversubscribed,
        avoiding the OOM kills that naive UVM caching suffers;
    a \emph{framework-guided eviction} strategy
        that uses \texttt{cudaMemDiscardBatchAsync} to mark partially-free
        cached blocks as evictable without copying their data back to the host;
    and \emph{prefetching with location hints}
        (\texttt{cudaPrefetchAsync} + \texttt{cudaMemAdvise})
        to hide GPU page-fault latency on reuse.
We implement {\uvmname} on PyTorch~2.8 with CUDA~13
    and evaluate Hugging~Face OPT-13B inference
    on A100-PCIe, H100-PCIe, and GH200 NVLink-c2c systems.

When memory is not oversubscribed, {\uvmname} matches PyTorch's native
    \texttt{cudaMalloc}-based caching allocator within $4$--$20\%$.
In the low-oversubscription regime, where the native allocator crashes with OOM,
    {\uvmname} is the fastest option,
    delivering up to $45\times$--$58\times$ speedup over naive UVM
    and $8\times$--$35\times$ over Hugging~Face's KV-cache offloading.
Under high oversubscription {\uvmname} still avoids OOM,
    but narrows to $1.6\times$--$1.7\times$ over default UVM
    and falls behind HF offloading by $\sim\!3\times$---a gap we are actively closing
    via the thrashing-mitigation mechanisms described above.
\end{abstract}
